% Responsable: Fernando Hernández
\section{Introducción y público objetivo}

En el Laboratorio 1 analizamos el ingreso mensual de viajeros internacionales a
Guatemala entre enero de 2009 y junio de 2026: 210 meses de la Base de Migración
de INGUAT/IGM. Construimos siete series de visitantes (turistas más
excursionistas): el total del país, las tres regiones de origen con más viajeros
(América del Centro, América del Norte y Europa) y los tres puntos de entrada con
más flujo (Aeropuerto La Aurora, Valle Nuevo y San Cristóbal). Sobre ellas hicimos
descomposiciones, pruebas de estacionariedad, modelos ARIMA, Holt-Winters,
suavizamiento exponencial y Prophet, y medimos el error de cada modelo.

Ese trabajo responde preguntas técnicas. La infografía tiene que responder una
sola pregunta de negocio:

\begin{center}
\itshape ¿Qué nos enseñan 17 años de datos migratorios sobre el turismo en
Guatemala y qué debemos anticipar hacia el futuro?
\end{center}

\subsection{Para quién es}

El lector es un directivo de INGUAT. Conoce el sector turístico, decide sobre
presupuesto de promoción, campañas y coordinación con Migración, pero no es
estadístico. De eso salen tres condiciones de diseño que se aplican en toda la
pieza:

\begin{enumerate}[leftmargin=*]
  \item \textbf{Cada gráfico tiene que tener una conclusión escrita.} El título de
  cada visualización es la conclusión (``Diciembre recibe 35\% más visitantes que
  un mes típico'') y no la descripción (``Visitantes por mes'').
  \item \textbf{Nada de vocabulario técnico sin traducir.} ``Fuerza estacional''
  se lee como ``qué tan predecible es el calendario de la serie''; ``RMSE'' como
  ``error típico del pronóstico, en visitantes por mes''; ``intervalo de
  confianza'' como ``rango probable'', rotulado ``Estimación''.
  \item \textbf{Las cifras se dan en unidades que el lector usa.} Visitantes,
  millones de visitantes por año y porcentajes respecto a 2019, el último año
  normal antes de la pandemia, en lugar de logaritmos, diferencias o índices sin
  referencia.
\end{enumerate}

\subsection{Del laboratorio a la infografía}

No todo lo que hicimos en el Laboratorio 1 entra en la infografía. Usamos este
criterio: entra lo que cambia una decisión. Las pruebas ADF y las funciones de
autocorrelación justifican los modelos, pero el directivo no decide nada con
ellas, así que se quedan fuera. La tendencia, el calendario de picos, el golpe de
la pandemia, qué modelo pronostica mejor y cuánto se equivoca sí cambian
decisiones, y cada uno tiene un bloque.

La medida es siempre la misma: visitantes por mes. Es la única agregación
comparable en los 210 meses, porque la base cambia de metodología a mitad del
periodo. Usar una sola unidad en toda la página evita que el lector compare
cifras que no son comparables.

\subsection{Cómo está organizado este documento}

La sección~\ref{sec:recorrido} explica el orden de los bloques y el recorrido de
lectura, con el boceto de la página. Las secciones 3 a 5 justifican, bloque por
bloque, el tipo de gráfico elegido según el tipo de dato, la paleta y la
tipografía. La sección 6 presenta el cuadrante comparativo con las cifras que
lo respaldan.
